% !TEX root = main.tex
% 名称。策略名称为暂定名称，仅在此处统一修改。
\newcommand{\method}{TurnAhead\xspace}
\newcommand{\spork}{SPORK\xspace}
\newcommand{\stratLaunch}{锚定前瞻\xspace}
\newcommand{\stratCommit}{精确输入提交\xspace}
\newcommand{\stratRecover}{排空后恢复\xspace}

% 智能体循环的记号（与第 2 节共同定稿）。
\newcommand{\hist}{H}                % 已提交历史
\newcommand{\ctx}[1]{X_{#1}}         % 第 #1 轮的模型输入
\newcommand{\render}{F}              % 历史策略／提示词渲染器
\newcommand{\rsn}[1]{r_{#1}}         % 推理文本
\newcommand{\vis}[1]{v_{#1}}         % 可见的非推理文本
\newcommand{\act}[1]{a_{#1}}         % 动作（工具调用）
\newcommand{\obs}[1]{o_{#1}}         % 观测
\newcommand{\env}[1]{S_{#1}}         % 环境状态
\newcommand{\anchor}{\mathcal{A}}    % 锚点
\newcommand{\slack}{\sigma}          % 锚点就绪后的剩余时间
\newcommand{\contract}{I}            % 提交契约

% 定理类环境。
\theoremstyle{plain}
\newtheorem{proposition}{命题}
\newtheorem{lemma}{引理}
\newtheorem{corollary}{推论}
\theoremstyle{definition}
\newtheorem{definition}{定义}
\newtheorem{assumption}{假设}
\theoremstyle{remark}
\newtheorem{remark}{注}
\crefname{section}{节}{节}
\Crefname{section}{节}{节}
\crefformat{section}{#2第~#1~节#3}
\Crefformat{section}{#2第~#1~节#3}
\crefname{subsection}{节}{节}
\Crefname{subsection}{节}{节}
\crefformat{subsection}{#2第~#1~节#3}
\Crefformat{subsection}{#2第~#1~节#3}
\crefname{appendix}{附录}{附录}
\Crefname{appendix}{附录}{附录}
\crefname{equation}{式}{式}
\Crefname{equation}{式}{式}
\crefname{figure}{图}{图}
\Crefname{figure}{图}{图}
\crefname{table}{表}{表}
\Crefname{table}{表}{表}
\crefname{algorithm}{算法}{算法}
\Crefname{algorithm}{算法}{算法}
\crefname{proposition}{命题}{命题}
\Crefname{proposition}{命题}{命题}
\crefname{lemma}{引理}{引理}
\Crefname{lemma}{引理}{引理}
\crefname{corollary}{推论}{推论}
\Crefname{corollary}{推论}{推论}
\crefname{definition}{定义}{定义}
\Crefname{definition}{定义}{定义}
\crefname{assumption}{假设}{假设}
\Crefname{assumption}{假设}{假设}
\AtBeginDocument{
\renewcommand{\crefpairconjunction}{和}
\renewcommand{\crefmiddleconjunction}{、}
\renewcommand{\creflastconjunction}{和}
\renewcommand{\crefrangeconjunction}{至}
}
\floatname{algorithm}{算法}
% 中文过程名使用常规字形，避免切换到不支持中文的小型大写字形。
\renewcommand{\textproc}[1]{\textnormal{#1}}
\algrenewcommand{\algorithmicprocedure}{过程}
\algrenewcommand{\algorithmicif}{若}
\algrenewcommand{\algorithmicthen}{则}
\algrenewcommand{\algorithmicelse}{否则}
\algrenewcommand{\algorithmicend}{结束}
\algrenewcommand{\algorithmicreturn}{返回}

% 草稿辅助标记；投稿前删除所有用例。
\newcommand{\planned}[1]{{\color{gray}\itshape 【计划补充】#1}}
\newcommand{\todo}[1]{{\color{red}【待办：#1】}}
